% ID: SEL-D0059-OA-E02
% ESTADO: cotejada. Segunda pasada visual y borrador compilado revisados.
% Fuente: Selectividad/sel_2011_matematicas_aplicadas/222-EXAMEN 2.pdf
% SHA-256: 998f20402054ca6d0454da83ffc8e1ff6951a576b48836a32e81e524c5985214
% Página: 1; opción A; ejercicio original 2.

Un banco lanza al mercado un plan de inversión cuya rentabilidad $R(x)$, en miles de euros, viene dada en función de la cantidad, $x$, que se invierte, también en miles de euros, por la siguiente expresión:
\[
R(x)=-0.001x^2+0.4x+3.5,\ \text{con }x\geq10.
\]
\begin{enumerate}[label=\alph*)]
\item \textbf{(0.5 puntos)} Calcule la rentabilidad para una inversión de 100000 euros.
\item \textbf{(1.5 puntos)} Deduzca y razone qué cantidad habría que invertir para obtener la máxima rentabilidad.
\item \textbf{(0.5 puntos)} \textquestiondown{}Qué rentabilidad máxima se obtendría?
\end{enumerate}
